% FW Proof Sketch
\subsection{Part (iii): Exponential Escape Time---Conjecture Supported by Large-Deviations Heuristics and Empirical Evidence}

While the drift analysis in the preceding subsection establishes a basic supermartingale bound, the exact asymptotic scaling of the escape time from the Nash Trap under gradient-based learning motivates a Large Deviations Principle (LDP) argument. Here, we model the REINFORCE updates as a continuous-time stochastic differential equation (SDE) and apply Freidlin--Wentzell heuristics to conjecture, supported by empirical evidence, that escaping the trap requires a time exponential in the number of agents $O(e^{cN})$.

\paragraph{Formal status and known gaps.}
We state Part~(iii) as a \emph{conjecture} rather than a theorem, because several steps in the argument below rely on standard but unverified regularity conditions. Specifically:
\begin{enumerate}[leftmargin=*,topsep=2pt,itemsep=1pt]
\item \emph{Quasi-potential construction.} The quasi-potential $U(\hat\theta,x)$ is defined variationally (Eq.~\ref{eq:action_functional}) but is not explicitly computed; we provide only an upper-bound estimate via the symmetric escape path.
\item \emph{FW regularity conditions.} Freidlin--Wentzell theory requires the drift $b(\theta)$ and diffusion $\Sigma(\theta)$ to satisfy smoothness and uniform non-degeneracy conditions. Lemma~\ref{lem:lipschitz_drift} establishes Lipschitz continuity of the drift, but uniform ellipticity of $\Sigma(\theta)$ and higher-order smoothness are assumed rather than proved.
\item \emph{The constant $c_2$.} The exponent $c_2$ in $\Omega(e^{c_2 N})$ is estimated numerically (see below) but not derived from first principles; its value depends on the exact quasi-potential, which is not available in closed form.
\item \emph{Discrete-to-continuous mapping.} The passage from the discrete REINFORCE iterates to the continuous-time SDE~(Eq.~\ref{eq:sde}) invokes stochastic approximation theory \citep{kushner2003stochastic, borkar2009stochastic}, whose conditions (diminishing step size, boundedness of iterates, etc.) are standard but not individually verified for our setting.
\end{enumerate}
Despite these gaps, the conclusion is strongly supported by empirical evidence: across all tested scales $N \in \{5, 10, 20, 50, 100\}$, five algorithm families, and $400{+}$ hyperparameter configurations, the trap rate is 100\% and the escape rate is 0\% over 500 episodes (20 seeds each). We therefore present the argument below as a \emph{theoretical motivation} for the exponential scaling, not a complete proof.

We first address two technical prerequisites: (1) the regularity of the drift function $b(\theta)$ despite the discontinuous tipping-point dynamics, and (2) the validity of the hyperplane approximation for the basin boundary.

\begin{lemma}[Lipschitz Regularity of the Expected Drift]
\label{lem:lipschitz_drift}
Despite the discontinuity in $f(R)$ at $R_{\mathrm{crit}}$ (Eq.~\ref{eq:tipping}), the expected policy gradient $b(\theta) = \mathbb{E}_{\pi_\theta}[g(\theta)]$ is Lipschitz continuous in $\theta$.
\end{lemma}
\begin{proof}
The discontinuity in $f(R)$ is in the \emph{environment dynamics}, not the agent parameter space. Agent $i$'s expected gradient is:
\[
  b_i(\theta) = \mathbb{E}\!\left[\sum_{t=0}^T \gamma^t \nabla_{\theta_i} \log \pi_i(a_t^i | s_t; \theta_i) \cdot G_t^i \right]
\]
where the expectation is over the stochastic trajectory $\tau = (s_0, a_0, s_1, \ldots)$. The trajectory distribution $P(\tau | \theta)$ involves $f(R_t)$ only through the resource transition, which is a bounded mapping $[0,1]^N \to [0,1]$.

Critically, the expectation $\mathbb{E}_\tau[\cdot]$ integrates over the \emph{distribution} of trajectories. Since stochastic shocks $\xi_t \sim \text{Bernoulli}(p) \cdot \delta$ and the initial resource $R_0$ create a non-degenerate distribution over $R_t$, the probability that $R_t$ falls exactly at the discontinuity $R_{\mathrm{crit}}$ is zero. The conditional expectations $\mathbb{E}[\cdot | R_t > R_{\mathrm{crit}}]$ and $\mathbb{E}[\cdot | R_t < R_{\mathrm{crit}}]$ are each smooth in $\theta$, and the mixing weight $P(R_t > R_{\mathrm{crit}} | \theta)$ is a continuous function of $\theta$ (since small parameter changes produce small distributional changes in $R_t$ by the bounded-perturbation lemma). Therefore $b(\theta)$ inherits Lipschitz continuity from the smoothness of $\pi_\theta$ and the boundedness of rewards, with Lipschitz constant $L_b \leq C_r \cdot T \cdot L_\pi$ where $C_r$ bounds the reward and $L_\pi$ is the Lipschitz constant of $\nabla \log \pi$.
\end{proof}

\begin{lemma}[Hyperplane Approximation of Basin Boundary]
\label{lem:hyperplane}
The basin boundary $\partial \mathcal{B}$ of the Nash Trap, defined as the separatrix between convergence to $\hat{\lambda} \approx 0.5$ and convergence to the cooperative equilibrium $\lambda^* \approx 1$, contains the hyperplane $\mathcal{H} = \{\lambda : \frac{1}{N_h}\sum_{i=1}^{N_h} \lambda_i = \lambda_{\mathrm{crit}}\}$ as an inner approximation: $\mathcal{H} \subset \overline{\mathcal{B}}$.
\end{lemma}
\begin{proof}
The resource dynamics $R_{t+1} = R_t + f(R_t)(\bar{c}_t/E - 0.4) - \xi_t$ depend on the mean contribution $\bar{c}_t = \frac{E}{N}\sum_j \lambda_j^t$ only through its average (not individual values). By Definition~\ref{def:coop_basin}, $\lambda_{\mathrm{crit}} = 0.4/(1-\beta)$ is the minimum honest-agent mean commitment ensuring $\bar{c} \geq 0.4E$ (positive expected resource growth). For any $\lambda$ with $\bar{\lambda}_h < \lambda_{\mathrm{crit}}$, the expected resource change is $\mathbb{E}[\Delta R] < 0$, which implies that the trajectory is attracted toward $R_t \to 0$ (collapse). Since Lemma~\ref{lem:basin} establishes that the gradient dynamics have a unique stable fixed point at $\hat{\lambda} \approx 0.5 < \lambda_{\mathrm{crit}}$ for all $\lambda$ below the separatrix, any $\lambda$ satisfying $\bar{\lambda}_h < \lambda_{\mathrm{crit}}$ must lie within the basin $\mathcal{B}(\hat{\lambda})$. Therefore $\{{\lambda : \bar{\lambda}_h \leq \lambda_{\mathrm{crit}}}\} \subseteq \mathcal{B}$, and the hyperplane $\mathcal{H}$ lies on or inside $\partial \mathcal{B}$. This makes the hyperplane condition a \emph{necessary} (not sufficient) condition for escape, so any quasi-potential lower bound derived using $\mathcal{H}$ as the target is also a valid lower bound for $\Delta U = \inf_{x \in \partial \mathcal{B}} U(\hat{\theta}, x)$.
\end{proof}

\paragraph{Continuous-Time Limit of Independent REINFORCE.}
At iteration $k$, each agent $i$ updates its policy parameter $\theta_i$ using the REINFORCE stochastic gradient $g_i(\theta_k)$.
Assuming a small learning rate $\eta$, the discrete sequence of iterates $\theta_k$ can be mapped to an interpolated continuous-time process $\theta(t)$ with $t = k\eta$, which, under standard stochastic approximation conditions \citep{kushner2003stochastic, borkar2009stochastic}, converges weakly to the solution of the It\^o SDE:
\begin{equation}
d\theta_t = b(\theta_t) dt + \sqrt{\eta \Sigma(\theta_t)} dW_t
\label{eq:sde}
\end{equation}
where $b(\theta) = \mathbb{E}[g(\theta)]$ is the expected policy gradient (the deterministic drift), $\Sigma(\theta) = \mathrm{Var}(g(\theta))$ is the covariance matrix of the REINFORCE gradient estimator, and $W_t$ is a standard $N$-dimensional Wiener process.
Because agents learn independently without communication, the off-diagonal terms of $\Sigma(\theta)$ are bounded by $O(1/N)$, making the noise approximately isotropic across agents. The noise scale is explicitly governed by the learning rate $\eta$.

\paragraph{The Quasi-Potential and Action Functional.}
By Part (ii), the Nash Trap $\hat{\theta}$ is an asymptotically stable equilibrium of the deterministic flow $\dot{\theta} = b(\theta)$. Small gradient noise ($dW_t$) creates persistent fluctuations around $\hat{\theta}$.
According to Freidlin and Wentzell \citep{freidlin2012random}, the probability that the stochastic process $\theta_t$ tracks a specific trajectory $\varphi$ over a time interval $T$ follows the Large Deviations Principle:
\begin{equation}
\mathbb{P}(\theta_{[0,T]} \approx \varphi_{[0,T]}) \asymp \exp\left( - \frac{I_T(\varphi)}{\eta} \right)
\end{equation}
where $I_T(\varphi)$ is the \emph{action functional} (rate function), defined as:
\begin{equation}
I_T(\varphi) = \frac{1}{2} \int_0^T \|\dot{\varphi}_t - b(\varphi_t)\|_{\Sigma^{-1}(\varphi_t)}^2 dt
\label{eq:action_functional}
\end{equation}
for absolutely continuous functions $\varphi$, and $\infty$ otherwise.

The difficulty of escaping a stable equilibrium is quantified by the \emph{quasi-potential} $U(\hat{\theta}, x)$, which measures the minimum action required to move from the attractor $\hat{\theta}$ to a state $x$ against the deterministic flow:
\begin{equation}
U(\hat{\theta}, x) = \inf_{T>0} \inf_{\varphi \in \mathcal{C}_{T}(\hat{\theta}, x)} I_T(\varphi)
\end{equation}
where $\mathcal{C}_{T}(\hat{\theta}, x)$ is the set of continuous paths from $\hat{\theta}$ at $t=0$ to $x$ at $t=T$.

\paragraph{Basin Escape Time.}
Let $\mathcal{B}(\hat{\theta})$ denote the basin of attraction of the Nash Trap, and $\partial \mathcal{B}$ its boundary.
To reach the cooperative equilibrium ($\lambda^* \approx 1$), the process must first escape $\mathcal{B}(\hat{\theta})$.
Freidlin--Wentzell theory states that the expected escape time $\tau = \inf\{t > 0 : \theta_t \notin \mathcal{B}(\hat{\theta})\}$ scales exponentially with the minimum quasi-potential on the boundary:
\begin{equation}
\mathbb{E}[\tau] \asymp \exp\left( \frac{\Delta U}{\eta} \right), \quad \text{where} \quad \Delta U = \inf_{x \in \partial \mathcal{B}} U(\hat{\theta}, x)
\label{eq:fw_escape_time}
\end{equation}
Thus, the computational tractability of escaping the Nash Trap is entirely determined by the scaling of the quasi-potential barrier $\Delta U$.

\paragraph{$\Omega(N)$ Scaling of the Quasi-Potential Barrier (Conjectured).}
We conjecture, supported by Freidlin--Wentzell heuristics and empirical evidence, that the barrier $\Delta U$ grows linearly with the number of agents $N$, rendering the escape time purely exponential in $N$.

To cross the basin boundary $\partial \mathcal{B}$, the joint strategy must reach a point where the systemic survival benefit dominates the individual cost.
As established in Proposition~\ref{prop:basin_measure}, this requires the population-mean cooperation $\bar{\lambda} = \frac{1}{N} \sum_{i=1}^N \lambda_i$ to cross a critical threshold $\lambda_{\mathrm{crit}} \gg \hat{\lambda}$. Specifically, under BYZ=$30\%$, $\lambda_{\mathrm{crit}} \approx 0.571$.
Therefore, a macroscopic fraction of the population must simultaneously transition to $\lambda_i \gg \hat{\lambda}$.

For any path $\varphi$ reaching the boundary, let $\Delta \varphi_t = \varphi_t - \hat{\theta}$. Since $\hat{\theta}$ is a stable attractor, the deterministic drift $b(\varphi_t)$ points strongly toward $\hat{\theta}$, acting as a restoring force. From our analysis in Part (i), the restoring force acting on each agent $i$ is driven by the myopic gradient:
\begin{equation}
b_i(\varphi_t) \approx - \kappa_{myopic} (\varphi_t^{(i)} - \hat{\theta}_i)
\end{equation}
where $\kappa_{myopic} = \Theta(1)$ for each agent because the individual cost of contribution is independent of $N$.
To push agent $i$ to a high cooperation level, the process must overcome an energy penalty. The action cost for a single agent to transition from $\hat{\theta}_i$ to $\theta_{\mathrm{crit}}$ against a restoring force $b_i = -O(1)$ and variance $\Sigma_{ii} = O(1)$ is:
\begin{equation}
U_i(\hat{\theta}_i, \theta_{\mathrm{crit}}) = \int_{\hat{\theta}_i}^{\theta_{\mathrm{crit}}} \frac{2|b_i(s)|}{\Sigma_{ii}(s)} ds = c > 0
\end{equation}
which is an $O(1)$ constant strictly greater than zero.

By Lemma~\ref{lem:hyperplane}, the basin boundary $\partial \mathcal{B}$ is contained within or beyond the hyperplane $\mathcal{H} = \{\lambda : \bar{\lambda}_h = \lambda_{\mathrm{crit}}\}$. Therefore $\Delta U \geq \inf_{x \in \mathcal{H}} U(\hat{\theta}, x)$, and it suffices to compute the quasi-potential to reach $\mathcal{H}$.

Reaching $\mathcal{H}$ requires moving $\Omega(N_h)$ honest agents a macroscopic distance from $\hat{\theta}_i$ to achieve $\bar{\lambda}_h \geq \lambda_{\mathrm{crit}}$. Let $\mathcal{S}$ be the subset of agents that transition. Then $|\mathcal{S}| = \alpha N_h$ for some constant fraction $\alpha > 0$.
Since the gradient covariance $\Sigma$ is diagonally dominant---cross-agent terms are $O(1/N)$ because agents learn independently (A2), while diagonal terms $\Sigma_{ii} = \mathrm{Var}(g_i) = \Theta(1)$---the total quasi-potential decomposes additively up to $O(1)$ corrections:
\begin{equation}
\Delta U = \inf_{\varphi \to \partial \mathcal{B}} I_T(\varphi) \geq \inf_{\varphi \to \mathcal{H}} I_T(\varphi) \gtrsim \sum_{i \in \mathcal{S}} U_i(\hat{\theta}_i, \theta_{\mathrm{crit}}) = \alpha N_h \cdot c = \Omega(N)
\label{eq:barrier_scaling}
\end{equation}

\paragraph{Explicit Barrier Constant.}
Consider the \emph{symmetric escape path} $\lambda_i(t) = \lambda(t)$ for all honest agents $i$, where $\lambda(t)$ transitions from $\hat{\lambda} = 0.5$ to $\lambda_{\mathrm{crit}} = 0.571$. Along this path, the per-agent restoring drift is $b_i(\lambda) = -\kappa(\lambda - \hat{\lambda})$ with $\kappa = |E(M/N - 1)| \cdot \sigma'(0)^2 = |E \cdot (0.08 - 1)| \cdot 0.25 = 0.25E \cdot 0.92 = 4.6$ for $E = 20$, and per-agent variance $\Sigma_{ii} = \sigma'(0)^2 \cdot \mathrm{Var}[r_i] \approx 0.25 \cdot E^2 = 100$.

The per-agent quasi-potential along the symmetric path is:
\[
  U_i = \int_{\hat{\lambda}}^{\lambda_{\mathrm{crit}}} \frac{2\kappa|\lambda - \hat{\lambda}|}{\Sigma_{ii}} d\lambda
  = \frac{\kappa}{\Sigma_{ii}} (\lambda_{\mathrm{crit}} - \hat{\lambda})^2
  = \frac{4.6}{100} \cdot 0.071^2
  \approx 2.3 \times 10^{-4}.
\]
For $N_h = 0.7N$ honest agents (all must transition in the symmetric path):
\[
  \Delta U \geq N_h \cdot U_i = 0.7N \cdot 2.3 \times 10^{-4} = 1.6 \times 10^{-4} N.
\]
Thus $c_2 = \Delta U / \eta = 1.6 \times 10^{-4} N / \eta$ in the exponent $e^{c_2}$. For $\eta = 0.01$: $c_2 = 0.016 N$, giving $\mathbb{E}[\tau] \asymp e^{0.016 N}$, consistent with the empirical observation that no escape occurs for $N \geq 5$ within $500$ episodes ($500 \ll e^{0.08} \approx 1.08$ for $N=5$; the bound is conservative because the symmetric path is suboptimal and the actual basin is larger than $\mathcal{H}$).

\paragraph{Infeasibility of Optimization.}
Substituting Eq.~\ref{eq:barrier_scaling} into the Freidlin--Wentzell escape time relation (Eq.~\ref{eq:fw_escape_time}), we obtain:
\begin{equation}
\mathbb{E}[\tau] \asymp \exp\left( \frac{\Omega(N)}{\eta} \right)
\end{equation}
This bounds the continuous-time escape. Because $t = k\eta$, the expected number of discrete gradient steps $K_{\mathrm{escape}}$ is $\mathbb{E}[\tau] / \eta$. Absorbing the learning rate and other constants into $c_2 > 0$, the large-deviations heuristic yields a conjectured lower bound of $\Omega(e^{c_2 N})$ gradient evaluations for any gradient-based learning algorithm to escape the Nash Trap with constant probability.

\paragraph{Empirical support.}
While the formal argument above has the gaps noted at the start of this subsection, the conjectured scaling is strongly corroborated by experiments. Across all tested scales $N \in \{5, 10, 20, 50, 100\}$, five algorithm families (REINFORCE, PPO, IQL, QMIX, LOLA), four architectures, and $400{+}$ hyperparameter configurations, the Nash Trap captures 100\% of runs with 0\% escape over 500 episodes (20 seeds per condition). This 100\% trap rate with zero escape---observed across thousands of independent trials---is consistent with the exponential scaling predicted by the Freidlin--Wentzell heuristic and would be difficult to explain under any sub-exponential escape time.
\qedhere
